\chapter{The Complete Source}
\label{app:source}

This is the healthy four-service example through chapter~\ref{ch:ownership}.
The deliberate failures are not included. Sessions owns the schedule,
Speakers owns speaker details and guarded email, Ratings contributes rating
fields, and Search owns a denormalized discovery projection.

Each heading below is a destination path from the example root. Create the
directories and save every file under that path. No project references
another project, even where two files have the same name. The schema exports
and router execution config are generated by the commands in
appendix~\ref{app:toolchain}; they are not additional source files to type.

Start from fresh exercise databases. The final state keeps the generated
loaders, nullable cross-service failure boundaries, the four-service
\code{graph.yaml} and token forwarding. Chapters~\ref{ch:satisfiability},
\ref{ch:breaking-changes} and~\ref{ch:ownership} use temporary inputs to
demonstrate failures; none replaces this baseline.

\section{Sessions}

\subsection*{\code{src/Sessions/ConferenceContext.cs}}
\begin{minted}{csharp}
using Microsoft.EntityFrameworkCore;

namespace Sessions;

/// <summary>
/// The conference program, in SQLite. Session carries SpeakerId as a
/// plain column, and since chapter 9 there is no speaker table here
/// for it to point at: the rows are the Speakers service's, and this
/// column is a foreign key across a seam that no database enforces.
/// </summary>
public sealed class ConferenceContext(
    DbContextOptions<ConferenceContext> options) : DbContext(options)
{
    public DbSet<Session> Sessions => Set<Session>();

    protected override void OnModelCreating(ModelBuilder model)
    {
        // SQLite stores no date type of its own and refuses to sort a
        // DateTimeOffset. Keep the offset out of the column and put it
        // back on the way in; every session here is in conference time.
        model.Entity<Session>()
            .Property(s => s.StartsAt)
            .HasConversion(
                v => v.UtcDateTime,
                v => new DateTimeOffset(v, TimeSpan.Zero));

        model.Entity<Session>().HasData(ConferenceData.Sessions);
    }
}
\end{minted}

\subsection*{\code{src/Sessions/ConferenceData.cs}}
\begin{minted}{csharp}
namespace Sessions;

/// <summary>
/// The conference program, seeded into this service's database. The
/// speakers that used to sit beside it moved to their own service in
/// chapter 9; what is left of them here is the SpeakerId on each row.
/// </summary>
public static class ConferenceData
{
    public static readonly IReadOnlyList<Session> Sessions =
    [
        new Session(
            1,
            "Schemas That Outlive Their Authors",
            "A five-year-old schema, and the decisions that shaped it.",
            1,
            new DateTimeOffset(2026, 9, 14, 9, 0, 0, TimeSpan.Zero),
            45),
        new Session(
            2,
            "Reading a Query Plan",
            null,
            1,
            new DateTimeOffset(2026, 9, 14, 10, 0, 0, TimeSpan.Zero),
            45),
        new Session(
            3,
            "Paging Without Offsets",
            "Why the page you get is not the page you asked for.",
            2,
            new DateTimeOffset(2026, 9, 14, 11, 0, 0, TimeSpan.Zero),
            30),
        new Session(
            4,
            "The Bank That Split Its Graph",
            "Eighteen months, four teams, and the unwritten parts.",
            3,
            new DateTimeOffset(2026, 9, 14, 14, 0, 0, TimeSpan.Zero),
            60)
    ];
}
\end{minted}

\subsection*{\code{src/Sessions/Mutation.cs}}
\begin{minted}{csharp}
using Microsoft.EntityFrameworkCore;

namespace Sessions;

[MutationType]
public static class Mutation
{
    /// <summary>
    /// Moves one session to a new slot, keeping its length.
    /// </summary>
    [Error<SessionNotFoundException>]
    [Error<SpeakerDoubleBookedException>]
    public static async Task<Session> RescheduleSessionAsync(
        [ID<Session>] int sessionId,
        DateTimeOffset startsAt,
        ConferenceContext db,
        CancellationToken ct)
    {
        var session = await db.Sessions
            .FirstOrDefaultAsync(s => s.Id == sessionId, ct)
            ?? throw new SessionNotFoundException(sessionId);

        var ends = startsAt.AddMinutes(session.DurationMinutes);

        // The overlap test is done in memory on purpose. StartsAt goes
        // through a value converter, so SQLite cannot do arithmetic on
        // that column, and one speaker's program is a handful of rows.
        var alsoTheirs = await db.Sessions
            .Where(s => s.SpeakerId == session.SpeakerId)
            .Where(s => s.Id != session.Id)
            .ToListAsync(ct);

        var clash = alsoTheirs.FirstOrDefault(
            s => s.StartsAt < ends
                && startsAt < s.StartsAt.AddMinutes(s.DurationMinutes));

        if (clash is not null)
        {
            throw new SpeakerDoubleBookedException(
                session.SpeakerId,
                clash.Id);
        }

        var moved = session with { StartsAt = startsAt };
        db.Entry(session).CurrentValues.SetValues(moved);
        await db.SaveChangesAsync(ct);

        return moved;
    }
}
\end{minted}

\subsection*{\code{src/Sessions/NullableNodesField.cs}}
\begin{minted}{csharp}
using HotChocolate.Configuration;
using HotChocolate.Types.Descriptors;
using HotChocolate.Types.Descriptors.Configurations;

namespace Sessions;

/// <summary>
/// Makes the generated nodes field nullable, so that a subgraph it
/// routes to being down costs the field rather than the response.
/// </summary>
// AddGlobalObjectIdentification generates node and nodes, and it
// generates them with different nullability: node is Node, which may
// hold a null, and nodes is [Node]!, which may not. The item type
// inside that list is nullable and it does not help, because the
// router resolves nodes as one batched step per owning subgraph and
// fails the whole field when one of them is unreachable. A non-null
// list with nowhere to put the failure propagates to its parent, and
// the parent of a root field is the response.
//
// Neither field is declared in any class here, so there is no
// signature to put a question mark on. This is the same problem
// ShareNodeFields beside it solves for a directive, and the same hook
// solves it: the field's Type is settable on the configuration.
internal sealed class NullableNodesField : TypeInterceptor
{
    public override void OnBeforeCompleteType(
        ITypeCompletionContext context,
        TypeSystemConfiguration configuration)
    {
        if (configuration is not ObjectTypeConfiguration type
            || type.Name != "Query")
        {
            return;
        }

        foreach (var field in type.Fields)
        {
            if (field.Name is "nodes")
            {
                field.Type = TypeReference.Parse(
                    "[Node]", TypeContext.Output);
            }
        }
    }
}
\end{minted}

\subsection*{\code{src/Sessions/Program.cs}}
\begin{minted}{csharp}
using ChilliCream.Nitro.App;
using HotChocolate.AspNetCore;
using Microsoft.EntityFrameworkCore;
using Sessions;

var builder = WebApplication.CreateBuilder(args);

// EF Core reports every command through the logging pipeline as well,
// at Information and with a duration attached. StatementLog below is
// this service's instrument, so silence the other one rather than have
// every statement printed twice.
builder.Logging.AddFilter(
    "Microsoft.EntityFrameworkCore.Database.Command",
    LogLevel.Warning);

builder.Services.AddDbContextFactory<ConferenceContext>(options =>
{
    options.UseSqlite("Data Source=conference.db");

    if (builder.Environment.IsDevelopment())
    {
        options.AddInterceptors(new StatementLog());
    }
});

builder
    .AddGraphQL()
    .AddSessionsTypes()
    .RegisterDbContextFactory<ConferenceContext>()
    .AddQueryContext()
    .AddPagingArguments()
    .AddGlobalObjectIdentification()
    .AddMutationConventions()
    .AddApolloFederation()
    .ModifyCostOptions(options => options.ApplyCostDefaults = false)
    .TryAddTypeInterceptor<ShareNodeFields>()
    .TryAddTypeInterceptor<NullableNodesField>();

var app = builder.Build();

using (var scope = app.Services.CreateScope())
{
    var factory = scope.ServiceProvider
        .GetRequiredService<IDbContextFactory<ConferenceContext>>();
    using var db = factory.CreateDbContext();
    db.Database.EnsureCreated();
}

app.MapGraphQL()
    .WithOptions(options =>
    {
        // Serve the Nitro build that shipped in the package. The
        // default, ServeMode.Latest, downloads the current one
        // from ChilliCream's CDN instead.
        options.Tool.ServeMode = ServeMode.Embedded;
    });

app.RunWithGraphQLCommands(args);
\end{minted}

\subsection*{\code{src/Sessions/Properties/launchSettings.json}}
\begin{minted}[fontsize=\footnotesize,breakanywhere]{json}
{
  "$schema": "https://json.schemastore.org/launchsettings.json",
  "profiles": {
    "http": {
      "commandName": "Project",
      "dotnetRunMessages": true,
      "launchBrowser": false,
      "applicationUrl": "http://localhost:5001",
      "environmentVariables": {
        "ASPNETCORE_ENVIRONMENT": "Development"
      }
    }
  }
}
\end{minted}

\subsection*{\code{src/Sessions/Query.cs}}
\begin{minted}{csharp}
using GreenDonut.Data;
using Microsoft.EntityFrameworkCore;

namespace Sessions;

[QueryType]
public static class Query
{
    /// <summary>
    /// Sessions on the schedule, earliest first.
    /// </summary>
    [UsePaging(DefaultPageSize = 2, MaxPageSize = 10)]
    public static IQueryable<Session> GetSessions(
        ConferenceContext db,
        QueryContext<Session> query)
        => db.Sessions
            .OrderBy(s => s.StartsAt)
            .With(query);

    /// <summary>
    /// One session by its identifier, or null if no session has it.
    /// </summary>
    [GraphQLDeprecated(
        "Ask for `node(id:)` instead. This field takes the database "
        + "key, which is only unique among sessions.")]
    public static async Task<Session?> GetSessionByIdAsync(
        int id,
        ConferenceContext db,
        QueryContext<Session> query,
        CancellationToken ct)
        => await db.Sessions
            .Where(s => s.Id == id)
            .With(query)
            .FirstOrDefaultAsync(ct);
}
\end{minted}

\subsection*{\code{src/Sessions/Session.cs}}
\begin{minted}{csharp}
namespace Sessions;

/// <summary>
/// A talk on the conference program.
/// </summary>
/// <param name="Id">Stable identifier for this session.</param>
/// <param name="Title">The line that goes in the program.</param>
/// <param name="Abstract">The longer text, if there is one.</param>
/// <param name="SpeakerId">Who gives it. Not part of the schema.</param>
/// <param name="StartsAt">When it starts, in conference time.</param>
/// <param name="DurationMinutes">How long the slot is.</param>
public sealed record Session(
    int Id,
    string Title,
    string? Abstract,
    [property: GraphQLIgnore] int SpeakerId,
    DateTimeOffset StartsAt,
    int DurationMinutes);
\end{minted}

\subsection*{\code{src/Sessions/SessionDataLoader.cs}}
\begin{minted}{csharp}
using Microsoft.EntityFrameworkCore;

namespace Sessions;

public static class SessionDataLoader
{
    /// <summary>
    /// Every session one batch of representations asked for, keyed by
    /// identifier. The router sends this service a list rather than
    /// one key at a time, so the loader behind the reference resolver
    /// is what decides whether that list costs one statement or one
    /// statement each.
    /// </summary>
    // No projection here, and that is a decision rather than an
    // oversight. A QueryContext<Session> can be injected and does
    // narrow the SELECT, but the selector it builds fills every
    // property nobody asked for with a default, and Id is a property
    // nobody asks for. This dictionary is keyed on Id. Chapter 10
    // measures both halves and takes the columns rather than the risk.
    [DataLoader]
    public static async Task<Dictionary<int, Session>>
        GetSessionByIdAsync(
            IReadOnlyList<int> ids,
            ConferenceContext db,
            CancellationToken ct)
        => await db.Sessions
            .Where(s => ids.Contains(s.Id))
            .ToDictionaryAsync(s => s.Id, ct);
}
\end{minted}

\subsection*{\code{src/Sessions/SessionNotFoundException.cs}}
\begin{minted}{csharp}
using HotChocolate.Types.Relay;

namespace Sessions;

/// <summary>
/// No session on the program carries the identifier that was asked for.
/// </summary>
public sealed class SessionNotFoundException(int sessionId)
    : Exception($"No session has id {sessionId}.")
{
    /// <summary>
    /// The identifier that matched nothing.
    /// </summary>
    [ID<Session>]
    public int SessionId { get; } = sessionId;
}
\end{minted}

\subsection*{\code{src/Sessions/SessionType.cs}}
\begin{minted}{csharp}
using GreenDonut.Data;
using HotChocolate.ApolloFederation.Resolvers;
using HotChocolate.ApolloFederation.Types;
using HotChocolate.Types.Relay;
using Microsoft.EntityFrameworkCore;

namespace Sessions;

[ObjectType<Session>]
[Key("id")]
[ReferenceResolver(
    EntityResolver = nameof(ResolveSessionReferenceAsync))]
public static partial class SessionType
{
    /// <summary>
    /// The person giving this session.
    /// </summary>
    // Since chapter 9 this service holds the identifier and nothing
    // else, so the whole of the work is wrapping the foreign key it
    // already has in the type the schema promises. No statement is
    // issued and nothing is checked: if no speaker carries this
    // number, the Speakers service is where that is discovered.
    public static Speaker? GetSpeaker(
        [Parent(nameof(Session.SpeakerId))] Session session)
        => new(session.SpeakerId);

    /// <summary>
    /// Loads one session by the identifier inside a global node id.
    /// </summary>
    [NodeResolver]
    public static async Task<Session?> GetSessionAsync(
        int id,
        ConferenceContext db,
        QueryContext<Session> query,
        CancellationToken ct)
        => await db.Sessions
            .Where(s => s.Id == id)
            .With(query)
            .FirstOrDefaultAsync(ct);

    /// <summary>
    /// Answers one representation from the entity route. Two things
    /// differ from the node resolver above. The key arrives as the
    /// undecoded node id string, so this method decodes it. And the
    /// router hands this route a whole list rather than one key at a
    /// time, so the lookup goes through a DataLoader where the one
    /// above projects instead; SessionDataLoader says why it cannot
    /// do both.
    /// </summary>
    [GraphQLIgnore]
    public static async Task<Session?> ResolveSessionReferenceAsync(
        string id,
        INodeIdSerializer serializer,
        ISessionByIdDataLoader sessionById,
        CancellationToken ct)
    {
        var nodeId = serializer.Parse(id, typeof(int));

        // Nothing has checked what is inside this string, and the
        // router is not the only caller the entity route has. A
        // Speaker's node id decodes to a perfectly good integer;
        // without this line it would answer with the session that
        // happens to carry it.
        if (nodeId.TypeName != nameof(Session))
        {
            return null;
        }

        return await sessionById.LoadAsync(
            (int)nodeId.InternalId!, ct);
    }
}
\end{minted}

\subsection*{\code{src/Sessions/Sessions.csproj}}
\begin{minted}[fontsize=\footnotesize,breakanywhere]{xml}
<Project Sdk="Microsoft.NET.Sdk.Web">

  <PropertyGroup>
    <TargetFramework>net10.0</TargetFramework>
    <Nullable>enable</Nullable>
    <ImplicitUsings>enable</ImplicitUsings>
    <GenerateDocumentationFile>true</GenerateDocumentationFile>
    <NoWarn>$(NoWarn);CS1591</NoWarn>
  </PropertyGroup>

  <ItemGroup>
    <PackageReference Include="HotChocolate.ApolloFederation" Version="16.6.1" />
    <PackageReference Include="HotChocolate.AspNetCore" Version="16.6.1" />
    <PackageReference Include="HotChocolate.AspNetCore.CommandLine" Version="16.6.1" />
    <PackageReference Include="HotChocolate.Data.EntityFramework" Version="16.6.1" />
    <PackageReference Include="HotChocolate.Types.Analyzers" Version="16.6.1" />
    <PackageReference Include="Microsoft.EntityFrameworkCore.Sqlite" Version="10.0.11" />
  </ItemGroup>

</Project>
\end{minted}

\subsection*{\code{src/Sessions/ShareNodeFields.cs}}
\begin{minted}{csharp}
using HotChocolate.Configuration;
using HotChocolate.Language;
using HotChocolate.Types.Descriptors.Configurations;

namespace Sessions;

/// <summary>
/// Applies @shareable to the two fields global object identification
/// generates, and to nothing else.
/// </summary>
internal sealed class ShareNodeFields : TypeInterceptor
{
    public override void OnBeforeCompleteType(
        ITypeCompletionContext context,
        TypeSystemConfiguration configuration)
    {
        if (configuration is not ObjectTypeConfiguration type
            || type.Name != "Query")
        {
            return;
        }

        foreach (var field in type.Fields)
        {
            if (field.Name is "node" or "nodes")
            {
                field.Directives.Add(
                    new DirectiveConfiguration(
                        new DirectiveNode("shareable")));
            }
        }
    }
}
\end{minted}

\subsection*{\code{src/Sessions/Speaker.cs}}
\begin{minted}{csharp}
using HotChocolate.Types.Relay;

namespace Sessions;

/// <summary>
/// A speaker, as far as this service is concerned.
/// </summary>
/// <param name="Id">Stable identifier for this speaker.</param>
// An identifier and nothing else. The name and the bio moved to the
// Speakers service in chapter 9, so there is no row behind this record
// and no table to load it from. It is worth declaring anyway, because
// a client asking a session for its speaker's name needs somewhere to
// ask, and this type is that somewhere.
public sealed record Speaker([property: ID<Speaker>] int Id);
\end{minted}

\subsection*{\code{src/Sessions/SpeakerDoubleBookedException.cs}}
\begin{minted}{csharp}
using HotChocolate.Types.Relay;

namespace Sessions;

/// <summary>
/// The new slot overlaps another session the same speaker gives.
/// </summary>
public sealed class SpeakerDoubleBookedException(
    int speakerId,
    int clashingSessionId)
    : Exception(
        $"Speaker {speakerId} already gives session "
        + $"{clashingSessionId} in that slot.")
{
    /// <summary>
    /// The speaker who cannot be in two rooms at once.
    /// </summary>
    [ID<Speaker>]
    public int SpeakerId { get; } = speakerId;

    /// <summary>
    /// The session already occupying the slot.
    /// </summary>
    [ID<Session>]
    public int ClashingSessionId { get; } = clashingSessionId;
}
\end{minted}

\subsection*{\code{src/Sessions/SpeakerType.cs}}
\begin{minted}{csharp}
using HotChocolate.ApolloFederation.Types;

namespace Sessions;

/// <summary>
/// The stub. This service names Speaker so that Session.speaker has a
/// type, declares which field identifies one, and then says it cannot
/// turn such a key back into anything: resolvable false is the whole
/// difference between this and the Speakers service's declaration of
/// the same type. Without it, composition would believe two subgraphs
/// can both answer for a speaker and route accordingly.
/// </summary>
[ObjectType<Speaker>]
[Key("id", false)]
public static partial class SpeakerType;
\end{minted}

\subsection*{\code{src/Sessions/StatementLog.cs}}
\begin{minted}{csharp}
using System.Data.Common;
using Microsoft.EntityFrameworkCore.Diagnostics;

namespace Sessions;

/// <summary>
/// Prints every statement the service sends to SQLite, numbered, so a
/// request's cost can be read off the console. The counter is
/// process-wide and never reset: start the service, send one request,
/// and the last number is what that request cost.
/// </summary>
public sealed class StatementLog : DbCommandInterceptor
{
    private static int _count;

    public override InterceptionResult<DbDataReader> ReaderExecuting(
        DbCommand command,
        CommandEventData eventData,
        InterceptionResult<DbDataReader> result)
    {
        Print(command);
        return result;
    }

    public override ValueTask<InterceptionResult<DbDataReader>>
        ReaderExecutingAsync(
            DbCommand command,
            CommandEventData eventData,
            InterceptionResult<DbDataReader> result,
            CancellationToken cancellationToken = default)
    {
        Print(command);
        return ValueTask.FromResult(result);
    }

    private static void Print(DbCommand command)
    {
        var n = Interlocked.Increment(ref _count);
        Console.WriteLine($"-- statement {n}");
        Console.WriteLine(command.CommandText);
    }
}
\end{minted}

\subsection*{\code{src/Sessions/appsettings.Development.json}}
\begin{minted}[fontsize=\footnotesize,breakanywhere]{json}
{
  "Logging": {
    "LogLevel": {
      "Default": "Information",
      "Microsoft.AspNetCore": "Warning"
    }
  }
}
\end{minted}

\subsection*{\code{src/Sessions/appsettings.json}}
\begin{minted}[fontsize=\footnotesize,breakanywhere]{json}
{
  "Logging": {
    "LogLevel": {
      "Default": "Information",
      "Microsoft.AspNetCore": "Warning"
    }
  },
  "AllowedHosts": "*"
}
\end{minted}
\section{Speakers}

\subsection*{\code{src/Speakers/NullableNodesField.cs}}
\begin{minted}{csharp}
using HotChocolate.Configuration;
using HotChocolate.Types.Descriptors;
using HotChocolate.Types.Descriptors.Configurations;

namespace Speakers;

/// <summary>
/// Makes the generated nodes field nullable, so that a subgraph it
/// routes to being down costs the field rather than the response.
/// Declared again here so both services expose the chosen nullable
/// boundary. The composer can merge different nullability into one
/// nullable field; matching the local contracts is this book's policy,
/// not a requirement that composition enforces.
/// </summary>
// AddGlobalObjectIdentification generates node and nodes, and it
// generates them with different nullability: node is Node, which may
// hold a null, and nodes is [Node]!, which may not. The item type
// inside that list is nullable and it does not help, because the
// router resolves nodes as one batched step per owning subgraph and
// fails the whole field when one of them is unreachable. A non-null
// list with nowhere to put the failure propagates to its parent, and
// the parent of a root field is the response.
//
// Neither field is declared in any class here, so there is no
// signature to put a question mark on. This is the same problem
// ShareNodeFields beside it solves for a directive, and the same hook
// solves it: the field's Type is settable on the configuration.
internal sealed class NullableNodesField : TypeInterceptor
{
    public override void OnBeforeCompleteType(
        ITypeCompletionContext context,
        TypeSystemConfiguration configuration)
    {
        if (configuration is not ObjectTypeConfiguration type
            || type.Name != "Query")
        {
            return;
        }

        foreach (var field in type.Fields)
        {
            if (field.Name is "nodes")
            {
                field.Type = TypeReference.Parse(
                    "[Node]", TypeContext.Output);
            }
        }
    }
}
\end{minted}

\subsection*{\code{src/Speakers/Program.cs}}
\begin{minted}{csharp}
using System.Text;
using ChilliCream.Nitro.App;
using HotChocolate.AspNetCore;
using Microsoft.EntityFrameworkCore;
using Microsoft.IdentityModel.Tokens;
using Speakers;

var builder = WebApplication.CreateBuilder(args);

// EF Core reports every command through the logging pipeline as well,
// at Information and with a duration attached. StatementLog below is
// this service's instrument, so silence the other one rather than have
// every statement printed twice.
builder.Logging.AddFilter(
    "Microsoft.EntityFrameworkCore.Database.Command",
    LogLevel.Warning);

builder.Services.AddDbContextFactory<SpeakerContext>(options =>
{
    options.UseSqlite("Data Source=speakers.db");

    if (builder.Environment.IsDevelopment())
    {
        options.AddInterceptors(new StatementLog());
    }
});

// A shared secret rather than a JWKS endpoint, because this book
// stands up no identity provider and a reader has to be able to mint
// a token and repeat every request in the chapter. The three
// Validate flags are off for the same reason and are the three a
// deployment turns on first. What is not negotiable is that the
// router validates the same token against the same key: a subgraph
// that trusts a header the router set, without checking it, has
// moved the guard back out of the service again.
builder.Services
    .AddAuthentication()
    .AddJwtBearer(options =>
    {
        options.TokenValidationParameters = new()
        {
            ValidateIssuer = false,
            ValidateAudience = false,
            ValidateLifetime = false,
            IssuerSigningKey = new SymmetricSecurityKey(
                Encoding.UTF8.GetBytes(SigningKey.Value))
        };
    });

builder.Services.AddAuthorization();

builder
    .AddGraphQL()
    // Without this the [Authorize] attributes compile, reach the
    // schema as @authorize, and are never executed.
    .AddAuthorization()
    .AddSpeakersTypes()
    .RegisterDbContextFactory<SpeakerContext>()
    .AddQueryContext()
    .AddPagingArguments()
    .AddGlobalObjectIdentification()
    .AddApolloFederation()
    .ModifyCostOptions(options => options.ApplyCostDefaults = false)
    .TryAddTypeInterceptor<ShareNodeFields>()
    .TryAddTypeInterceptor<NullableNodesField>();

var app = builder.Build();

using (var scope = app.Services.CreateScope())
{
    var factory = scope.ServiceProvider
        .GetRequiredService<IDbContextFactory<SpeakerContext>>();
    using var db = factory.CreateDbContext();
    db.Database.EnsureCreated();
}

app.UseAuthentication();
app.UseAuthorization();

app.MapGraphQL()
    .WithOptions(options =>
    {
        // Serve the Nitro build that shipped in the package. The
        // default, ServeMode.Latest, downloads the current one
        // from ChilliCream's CDN instead.
        options.Tool.ServeMode = ServeMode.Embedded;
    });

app.RunWithGraphQLCommands(args);
\end{minted}

\subsection*{\code{src/Speakers/Properties/launchSettings.json}}
\begin{minted}[fontsize=\footnotesize,breakanywhere]{json}
{
  "$schema": "https://json.schemastore.org/launchsettings.json",
  "profiles": {
    "http": {
      "commandName": "Project",
      "dotnetRunMessages": true,
      "launchBrowser": false,
      "applicationUrl": "http://localhost:5002",
      "environmentVariables": {
        "ASPNETCORE_ENVIRONMENT": "Development"
      }
    }
  }
}
\end{minted}

\subsection*{\code{src/Speakers/Query.cs}}
\begin{minted}{csharp}
using GreenDonut.Data;

namespace Speakers;

[QueryType]
public static class Query
{
    /// <summary>
    /// Everyone speaking at the conference, in program order.
    /// </summary>
    [UsePaging(DefaultPageSize = 2, MaxPageSize = 10)]
    public static IQueryable<Speaker> GetSpeakers(
        SpeakerContext db,
        QueryContext<Speaker> query)
        => db.Speakers
            .OrderBy(s => s.Name)
            .With(query);
}
\end{minted}

\subsection*{\code{src/Speakers/ShareNodeFields.cs}}
\begin{minted}{csharp}
using HotChocolate.Configuration;
using HotChocolate.Language;
using HotChocolate.Types.Descriptors.Configurations;

namespace Speakers;

/// <summary>
/// Applies @shareable to the two fields global object identification
/// generates, and to nothing else. This is chapter 7's file, declared
/// again in this service rather than shared with the Sessions one:
/// Sessions and Speakers both export node and nodes, so both need
/// this interceptor. Ratings and Search export neither. A project
/// reference between services that could live in separate repositories
/// would couple their builds.
/// </summary>
internal sealed class ShareNodeFields : TypeInterceptor
{
    public override void OnBeforeCompleteType(
        ITypeCompletionContext context,
        TypeSystemConfiguration configuration)
    {
        if (configuration is not ObjectTypeConfiguration type
            || type.Name != "Query")
        {
            return;
        }

        foreach (var field in type.Fields)
        {
            if (field.Name is "node" or "nodes")
            {
                field.Directives.Add(
                    new DirectiveConfiguration(
                        new DirectiveNode("shareable")));
            }
        }
    }
}
\end{minted}

\subsection*{\code{src/Speakers/SigningKey.cs}}
\begin{minted}{csharp}
namespace Speakers;

/// <summary>
/// The symmetric key this service validates tokens with. A real
/// deployment reads a key from configuration and validates issuer,
/// audience and lifetime; this one is fixed so that a reader can mint
/// a token and reproduce every request in the chapter.
/// </summary>
public static class SigningKey
{
    public const string Value =
        "splitting-the-graph-development-signing-key-32b";
}
\end{minted}

\subsection*{\code{src/Speakers/Speaker.cs}}
\begin{minted}{csharp}
using HotChocolate.ApolloFederation.Types;
using HotChocolate.Authorization;

namespace Speakers;

/// <summary>
/// A person giving one or more sessions. This service owns the rows;
/// the Sessions subgraph knows only that a speaker with this id exists.
/// </summary>
/// <param name="Id">Stable identifier for this speaker.</param>
/// <param name="Name">The name to print in the program.</param>
/// <param name="Bio">A sentence or two, if the speaker sent any.</param>
/// <param name="Email">Where the program committee writes. Not for
/// the printed program, and guarded on every route into it.</param>
public sealed record Speaker(
    int Id,
    string Name,
    string? Bio,
    [property: Authorize, Authenticated] string Email);
\end{minted}

\subsection*{\code{src/Speakers/SpeakerContext.cs}}
\begin{minted}{csharp}
using Microsoft.EntityFrameworkCore;

namespace Speakers;

/// <summary>
/// This service's own database, in its own file. Nothing here knows
/// that sessions exist: a speaker does not carry a list of them, and
/// no query in this service joins across that seam because there is
/// nothing on the other side of it to join to.
/// </summary>
public sealed class SpeakerContext(
    DbContextOptions<SpeakerContext> options) : DbContext(options)
{
    public DbSet<Speaker> Speakers => Set<Speaker>();

    protected override void OnModelCreating(ModelBuilder model)
        => model.Entity<Speaker>().HasData(SpeakerData.Speakers);
}
\end{minted}

\subsection*{\code{src/Speakers/SpeakerData.cs}}
\begin{minted}{csharp}
namespace Speakers;

/// <summary>
/// The speaker list, seeded into this service's own database. These are
/// the rows chapter 3 seeded into the Sessions service; they live here
/// now, and only here.
/// </summary>
public static class SpeakerData
{
    public static readonly IReadOnlyList<Speaker> Speakers =
    [
        new Speaker(
            1, "Ada Fischer", "Works on query planners.",
            "ada@example.test"),
        new Speaker(
            2, "Bruno Kaminski", null,
            "bruno@example.test"),
        new Speaker(
            3, "Chidi Okafor", "Runs platform engineering.",
            "chidi@example.test")
    ];
}
\end{minted}

\subsection*{\code{src/Speakers/SpeakerDataLoader.cs}}
\begin{minted}{csharp}
using Microsoft.EntityFrameworkCore;

namespace Speakers;

public static class SpeakerDataLoader
{
    /// <summary>
    /// Every speaker one batch asked for, keyed by identifier. The
    /// router sends this service a list of representations rather than
    /// one at a time, so the loader behind the reference resolver is
    /// what decides whether that list costs one statement or several.
    /// </summary>
    [DataLoader]
    public static async Task<Dictionary<int, Speaker>>
        GetSpeakerByIdAsync(
            IReadOnlyList<int> ids,
            SpeakerContext db,
            CancellationToken ct)
        => await db.Speakers
            .Where(s => ids.Contains(s.Id))
            .ToDictionaryAsync(s => s.Id, ct);
}
\end{minted}

\subsection*{\code{src/Speakers/SpeakerType.cs}}
\begin{minted}{csharp}
using HotChocolate.ApolloFederation.Resolvers;
using HotChocolate.ApolloFederation.Types;
using HotChocolate.Types.Relay;

namespace Speakers;

[ObjectType<Speaker>]
[Key("id")]
[ReferenceResolver(
    EntityResolver = nameof(ResolveSpeakerReferenceAsync))]
public static partial class SpeakerType
{
    /// <summary>
    /// Loads one speaker by the identifier inside a global node id.
    /// </summary>
    [NodeResolver]
    public static async Task<Speaker?> GetSpeakerAsync(
        int id,
        ISpeakerByIdDataLoader speakerById,
        CancellationToken ct)
        => await speakerById.LoadAsync(id, ct);

    /// <summary>
    /// Answers one representation from the entity route. This is the
    /// method the router reaches every time a client asks a session
    /// for its speaker's name, so it is the busiest resolver in the
    /// graph. The key arrives as the undecoded node id string, and the
    /// type name inside it is checked before the integer is spent.
    /// </summary>
    [GraphQLIgnore]
    public static async Task<Speaker?> ResolveSpeakerReferenceAsync(
        string id,
        INodeIdSerializer serializer,
        ISpeakerByIdDataLoader speakerById,
        CancellationToken ct)
    {
        var nodeId = serializer.Parse(id, typeof(int));

        if (nodeId.TypeName != nameof(Speaker))
        {
            return null;
        }

        return await speakerById.LoadAsync(
            (int)nodeId.InternalId!, ct);
    }
}
\end{minted}

\subsection*{\code{src/Speakers/Speakers.csproj}}
\begin{minted}[fontsize=\footnotesize,breakanywhere]{xml}
<Project Sdk="Microsoft.NET.Sdk.Web">

  <PropertyGroup>
    <TargetFramework>net10.0</TargetFramework>
    <Nullable>enable</Nullable>
    <ImplicitUsings>enable</ImplicitUsings>
    <GenerateDocumentationFile>true</GenerateDocumentationFile>
    <NoWarn>$(NoWarn);CS1591</NoWarn>
  </PropertyGroup>

  <ItemGroup>
    <PackageReference Include="HotChocolate.ApolloFederation" Version="16.6.1" />
    <PackageReference Include="HotChocolate.AspNetCore" Version="16.6.1" />
    <PackageReference Include="HotChocolate.AspNetCore.Authorization" Version="16.6.1" />
    <PackageReference Include="HotChocolate.AspNetCore.CommandLine" Version="16.6.1" />
    <PackageReference Include="HotChocolate.Data.EntityFramework" Version="16.6.1" />
    <PackageReference Include="HotChocolate.Types.Analyzers" Version="16.6.1" />
    <PackageReference Include="Microsoft.AspNetCore.Authentication.JwtBearer" Version="10.0.11" />
    <PackageReference Include="Microsoft.EntityFrameworkCore.Sqlite" Version="10.0.11" />
  </ItemGroup>

</Project>
\end{minted}

\subsection*{\code{src/Speakers/StatementLog.cs}}
\begin{minted}{csharp}
using System.Data.Common;
using Microsoft.EntityFrameworkCore.Diagnostics;

namespace Speakers;

/// <summary>
/// Prints every statement the service sends to SQLite, numbered, so a
/// request's cost can be read off the console. The counter is
/// process-wide and never reset: start the service, send one request,
/// and the last number is what that request cost.
/// </summary>
public sealed class StatementLog : DbCommandInterceptor
{
    private static int _count;

    public override InterceptionResult<DbDataReader> ReaderExecuting(
        DbCommand command,
        CommandEventData eventData,
        InterceptionResult<DbDataReader> result)
    {
        Print(command);
        return result;
    }

    public override ValueTask<InterceptionResult<DbDataReader>>
        ReaderExecutingAsync(
            DbCommand command,
            CommandEventData eventData,
            InterceptionResult<DbDataReader> result,
            CancellationToken cancellationToken = default)
    {
        Print(command);
        return ValueTask.FromResult(result);
    }

    private static void Print(DbCommand command)
    {
        var n = Interlocked.Increment(ref _count);
        Console.WriteLine($"-- statement {n}");
        Console.WriteLine(command.CommandText);
    }
}
\end{minted}

\subsection*{\code{src/Speakers/appsettings.Development.json}}
\begin{minted}[fontsize=\footnotesize,breakanywhere]{json}
{
  "Logging": {
    "LogLevel": {
      "Default": "Information",
      "Microsoft.AspNetCore": "Warning"
    }
  }
}
\end{minted}

\subsection*{\code{src/Speakers/appsettings.json}}
\begin{minted}[fontsize=\footnotesize,breakanywhere]{json}
{
  "Logging": {
    "LogLevel": {
      "Default": "Information",
      "Microsoft.AspNetCore": "Warning"
    }
  },
  "AllowedHosts": "*"
}
\end{minted}
\section{Ratings}

\subsection*{\code{src/Ratings/Program.cs}}
\begin{minted}{csharp}
using ChilliCream.Nitro.App;
using HotChocolate.AspNetCore;
using Microsoft.EntityFrameworkCore;
using Ratings;

var builder = WebApplication.CreateBuilder(args);

// EF Core reports every command through the logging pipeline as well,
// at Information and with a duration attached. StatementLog below is
// this service's instrument, so silence the other one rather than have
// every statement printed twice.
builder.Logging.AddFilter(
    "Microsoft.EntityFrameworkCore.Database.Command",
    LogLevel.Warning);

builder.Services.AddDbContextFactory<RatingContext>(options =>
{
    options.UseSqlite("Data Source=ratings.db");

    if (builder.Environment.IsDevelopment())
    {
        options.AddInterceptors(new StatementLog());
    }
});

// Two lines here differ from the other two services, and both follow
// from the same fact: this service contributes fields to a type it
// does not own and owns no node of its own.
//
// There is no AddGlobalObjectIdentification, because there cannot be.
// That call refuses to build a schema with no node type in it: "There
// is no object type implementing interface `Node`." So this service
// has no node or nodes field, and needs no interceptor to declare
// them shareable.
//
// But the keys it is handed across the seam are still global node ids,
// and the serializer that reads one is normally registered by the call
// it cannot make. AddDefaultNodeIdSerializer registers it on its own.
// Without this line the reference resolver is injected a null and the
// entity route answers Unexpected Execution Error.
builder
    .AddGraphQL()
    .AddRatingsTypes()
    .RegisterDbContextFactory<RatingContext>()
    .AddDefaultNodeIdSerializer()
    .AddApolloFederation()
    .ModifyCostOptions(options => options.ApplyCostDefaults = false);

var app = builder.Build();

using (var scope = app.Services.CreateScope())
{
    var factory = scope.ServiceProvider
        .GetRequiredService<IDbContextFactory<RatingContext>>();
    using var db = factory.CreateDbContext();
    db.Database.EnsureCreated();
}

app.MapGraphQL()
    .WithOptions(options =>
    {
        // Serve the Nitro build that shipped in the package. The
        // default, ServeMode.Latest, downloads the current one
        // from ChilliCream's CDN instead.
        options.Tool.ServeMode = ServeMode.Embedded;
    });

app.RunWithGraphQLCommands(args);
\end{minted}

\subsection*{\code{src/Ratings/Properties/launchSettings.json}}
\begin{minted}[fontsize=\footnotesize,breakanywhere]{json}
{
  "$schema": "https://json.schemastore.org/launchsettings.json",
  "profiles": {
    "http": {
      "commandName": "Project",
      "dotnetRunMessages": true,
      "launchBrowser": false,
      "applicationUrl": "http://localhost:5003",
      "environmentVariables": {
        "ASPNETCORE_ENVIRONMENT": "Development"
      }
    }
  }
}
\end{minted}

\subsection*{\code{src/Ratings/Query.cs}}
\begin{minted}{csharp}
using Microsoft.EntityFrameworkCore;

namespace Ratings;

[QueryType]
public static class Query
{
    /// <summary>
    /// How many scores this service holds, across every session.
    /// </summary>
    // The question mark is the whole of chapter 14. A count over this
    // service's own table is a value it can always produce, which is
    // what the nullability rule of chapter 4 asks of a non-null field.
    // Behind a router it is also a root field, and the nearest
    // position above a root field that is allowed to hold a null is
    // the response's own data entry. Declared non-null, this service
    // being stopped discards every other service's answer in the same
    // response along with it.
    public static async Task<int?> GetRatingCountAsync(
        RatingContext db,
        CancellationToken ct)
        => await db.Ratings.CountAsync(ct);
}
\end{minted}

\subsection*{\code{src/Ratings/Rating.cs}}
\begin{minted}{csharp}
namespace Ratings;

/// <summary>
/// One score somebody left for one session.
/// </summary>
/// <param name="Id">Stable identifier for this rating.</param>
/// <param name="SessionId">
/// Which session it is for. A number this service stores and cannot
/// check: the sessions are another service's rows.
/// </param>
/// <param name="Score">One to five.</param>
public sealed record Rating(int Id, int SessionId, int Score);
\end{minted}

\subsection*{\code{src/Ratings/RatingContext.cs}}
\begin{minted}{csharp}
using Microsoft.EntityFrameworkCore;

namespace Ratings;

/// <summary>
/// This service's own database. It holds ratings and nothing else: no
/// session table, no speaker table, and no foreign key that any
/// database could enforce, because the rows SessionId points at are
/// three miles away in another process.
/// </summary>
public sealed class RatingContext(
    DbContextOptions<RatingContext> options) : DbContext(options)
{
    public DbSet<Rating> Ratings => Set<Rating>();

    protected override void OnModelCreating(ModelBuilder model)
        => model.Entity<Rating>().HasData(RatingData.Ratings);
}
\end{minted}

\subsection*{\code{src/Ratings/RatingData.cs}}
\begin{minted}{csharp}
namespace Ratings;

/// <summary>
/// The scores this service was left with. Session 4 has none, which is
/// the case worth seeding deliberately: an average over no rows is not
/// zero, and the schema has to say which.
/// </summary>
public static class RatingData
{
    public static readonly IReadOnlyList<Rating> Ratings =
    [
        new Rating(1, 1, 5),
        new Rating(2, 1, 4),
        new Rating(3, 1, 3),
        new Rating(4, 2, 4),
        new Rating(5, 2, 5),
        new Rating(6, 3, 3)
    ];
}
\end{minted}

\subsection*{\code{src/Ratings/RatingDataLoader.cs}}
\begin{minted}{csharp}
using Microsoft.EntityFrameworkCore;

namespace Ratings;

public static class RatingDataLoader
{
    /// <summary>
    /// Every rating left for each session in one batch, grouped by
    /// session. A session nobody rated is absent from the dictionary.
    /// The field resolver treats a missing group or an empty group as
    /// no average and a count of zero.
    /// </summary>
    [DataLoader]
    public static async Task<Dictionary<int, Rating[]>>
        GetRatingsBySessionIdAsync(
            IReadOnlyList<int> sessionIds,
            RatingContext db,
            CancellationToken ct)
    {
        var rows = await db.Ratings
            .Where(r => sessionIds.Contains(r.SessionId))
            .ToListAsync(ct);

        return rows
            .GroupBy(r => r.SessionId)
            .ToDictionary(g => g.Key, g => g.ToArray());
    }
}
\end{minted}

\subsection*{\code{src/Ratings/Ratings.csproj}}
\begin{minted}[fontsize=\footnotesize,breakanywhere]{xml}
<Project Sdk="Microsoft.NET.Sdk.Web">

  <PropertyGroup>
    <TargetFramework>net10.0</TargetFramework>
    <Nullable>enable</Nullable>
    <ImplicitUsings>enable</ImplicitUsings>
    <GenerateDocumentationFile>true</GenerateDocumentationFile>
    <NoWarn>$(NoWarn);CS1591</NoWarn>
  </PropertyGroup>

  <ItemGroup>
    <PackageReference Include="HotChocolate.ApolloFederation" Version="16.6.1" />
    <PackageReference Include="HotChocolate.AspNetCore" Version="16.6.1" />
    <PackageReference Include="HotChocolate.AspNetCore.CommandLine" Version="16.6.1" />
    <PackageReference Include="HotChocolate.Data.EntityFramework" Version="16.6.1" />
    <PackageReference Include="HotChocolate.Types.Analyzers" Version="16.6.1" />
    <PackageReference Include="Microsoft.EntityFrameworkCore.Sqlite" Version="10.0.11" />
  </ItemGroup>

</Project>
\end{minted}

\subsection*{\code{src/Ratings/Session.cs}}
\begin{minted}{csharp}
using HotChocolate.Types.Relay;

namespace Ratings;

/// <summary>
/// A session, as far as this service is concerned.
/// </summary>
/// <param name="Id">Stable identifier for this session.</param>
/// <param name="Key">The number inside that identifier.</param>
/// <param name="Title">The line that goes in the program.</param>
// Declared here rather than referenced from the Sessions project. Two
// services that would be separate repositories owned by separate teams
// do not share a model assembly, and this book would be demonstrating
// the opposite of what it says if they did.
//
// Three fields where the Sessions service has one, and the split is
// the whole lesson of the seam. Id is the global node id, as a string,
// because that is what arrives and what has to go back. Key is the
// integer inside it, which is the only part this service's own tables
// know about, and it is hidden from the schema because no client has
// any use for it. Title is never stored here at all.
public sealed record Session(
    [property: ID] string Id,
    [property: GraphQLIgnore] int Key,
    string Title);
\end{minted}

\subsection*{\code{src/Ratings/SessionType.cs}}
\begin{minted}{csharp}
using HotChocolate.ApolloFederation.Resolvers;
using HotChocolate.ApolloFederation.Types;
using HotChocolate.Types.Relay;

namespace Ratings;

[ObjectType<Session>]
[Key("id")]
[ReferenceResolver(
    EntityResolver = nameof(ResolveSessionReference))]
public static partial class SessionType
{
    /// <summary>
    /// The line that goes in the program.
    /// </summary>
    // Owned by the Sessions service. Declared here only so that
    // feedbackUrl below has something to require, and resolved by
    // nobody: the router never asks this service for it.
    [External]
    public static string GetTitle([Parent] Session session)
        => session.Title;

    /// <summary>
    /// The mean score, or null if nobody has rated this session.
    /// </summary>
    public static async Task<double?> GetAverageScoreAsync(
        [Parent(nameof(Session.Key))] Session session,
        IRatingsBySessionIdDataLoader ratingsBySessionId,
        CancellationToken ct)
    {
        var ratings = await ratingsBySessionId
            .LoadAsync(session.Key, ct);

        return ratings is { Length: > 0 }
            ? ratings.Average(r => (double)r.Score)
            : null;
    }

    /// <summary>
    /// How many scores this session has been left.
    /// </summary>
    // Nullable for a different reason than the root field above. This
    // one is contributed to a type another service owns, so it rides
    // inside that service's list. Non-null here does not empty the
    // response; it empties the page of sessions the Sessions service
    // answered correctly, because a list of non-null items has nowhere
    // to put one null and goes null entirely.
    public static async Task<int?> GetRatingCountAsync(
        [Parent(nameof(Session.Key))] Session session,
        IRatingsBySessionIdDataLoader ratingsBySessionId,
        CancellationToken ct)
        => (await ratingsBySessionId.LoadAsync(session.Key, ct))
            ?.Length ?? 0;

    /// <summary>
    /// Where an attendee goes to leave a score for this session.
    /// </summary>
    // The url is built from the session's title, and the title belongs
    // to another service. @requires is what makes the router send it
    // along with the key, and [External] on GetTitle above is what
    // makes that declaration legal.
    //
    // Nullable on the same reasoning as ratingCount above, and the
    // @requires makes it the worse case rather than a better one: this
    // field cannot be answered unless the router has already collected
    // a title from a third service, so it has two ways to fail and
    // both of them land on somebody else's list.
    [Requires("title")]
    public static string? GetFeedbackUrl([Parent] Session session)
        => "https://feedback.example/2026/" + Slug(session.Title);

    /// <summary>
    /// Answers one representation from the entity route.
    /// </summary>
    // The key arrives as the undecoded node id string, exactly as it
    // does in the Sessions service, and for the same reason: nothing
    // on the entity route decodes one for you. Marking the property
    // [ID<Session>] does not do it either: it hands this method a zero
    // rather than raising anything, so the service answers with a
    // session that has no ratings and never says why.
    //
    // title is null on every request that does not ask for
    // feedbackUrl, because the router only sends a required field when
    // the field requiring it was asked for.
    [GraphQLIgnore]
    public static Session? ResolveSessionReference(
        string id,
        string? title,
        INodeIdSerializer serializer)
    {
        var nodeId = serializer.Parse(id, typeof(int));

        if (nodeId.TypeName != nameof(Session))
        {
            return null;
        }

        return new Session(
            id,
            (int)nodeId.InternalId!,
            title ?? string.Empty);
    }

    private static string Slug(string title)
        => string.Concat(
            title.ToLowerInvariant()
                .Select(c => char.IsLetterOrDigit(c) ? c : '-'));
}
\end{minted}

\subsection*{\code{src/Ratings/StatementLog.cs}}
\begin{minted}{csharp}
using System.Data.Common;
using Microsoft.EntityFrameworkCore.Diagnostics;

namespace Ratings;

/// <summary>
/// Prints every statement the service sends to SQLite, numbered, so a
/// request's cost can be read off the console. The counter is
/// process-wide and never reset: start the service, send one request,
/// and the last number is what that request cost.
/// </summary>
public sealed class StatementLog : DbCommandInterceptor
{
    private static int _count;

    public override InterceptionResult<DbDataReader> ReaderExecuting(
        DbCommand command,
        CommandEventData eventData,
        InterceptionResult<DbDataReader> result)
    {
        Print(command);
        return result;
    }

    public override ValueTask<InterceptionResult<DbDataReader>>
        ReaderExecutingAsync(
            DbCommand command,
            CommandEventData eventData,
            InterceptionResult<DbDataReader> result,
            CancellationToken cancellationToken = default)
    {
        Print(command);
        return ValueTask.FromResult(result);
    }

    private static void Print(DbCommand command)
    {
        var n = Interlocked.Increment(ref _count);
        Console.WriteLine($"-- statement {n}");
        Console.WriteLine(command.CommandText);
    }
}
\end{minted}

\subsection*{\code{src/Ratings/appsettings.Development.json}}
\begin{minted}[fontsize=\footnotesize,breakanywhere]{json}
{
  "Logging": {
    "LogLevel": {
      "Default": "Information",
      "Microsoft.AspNetCore": "Warning"
    }
  }
}
\end{minted}

\subsection*{\code{src/Ratings/appsettings.json}}
\begin{minted}[fontsize=\footnotesize,breakanywhere]{json}
{
  "Logging": {
    "LogLevel": {
      "Default": "Information",
      "Microsoft.AspNetCore": "Warning"
    }
  },
  "AllowedHosts": "*"
}
\end{minted}
\section{Search}

\subsection*{\code{src/Search/Program.cs}}
\begin{minted}{csharp}
using ChilliCream.Nitro.App;
using HotChocolate.AspNetCore;
using Microsoft.EntityFrameworkCore;
using Search;

var builder = WebApplication.CreateBuilder(args);

builder.Logging.AddFilter(
    "Microsoft.EntityFrameworkCore.Database.Command",
    LogLevel.Warning);

builder.Services.AddDbContextFactory<SearchContext>(options =>
{
    options.UseSqlite("Data Source=search.db");
    if (builder.Environment.IsDevelopment())
    {
        options.AddInterceptors(new StatementLog());
    }
});

builder
    .AddGraphQL()
    .AddSearchTypes()
    .RegisterDbContextFactory<SearchContext>()
    .AddPagingArguments()
    .AddDefaultNodeIdSerializer()
    .AddApolloFederation()
    .ModifyCostOptions(options => options.ApplyCostDefaults = false);

var app = builder.Build();

using (var scope = app.Services.CreateScope())
{
    var factory = scope.ServiceProvider
        .GetRequiredService<IDbContextFactory<SearchContext>>();
    using var db = factory.CreateDbContext();
    db.Database.EnsureCreated();
}

app.MapGraphQL().WithOptions(options =>
{
    options.Tool.ServeMode = ServeMode.Embedded;
});

app.RunWithGraphQLCommands(args);
\end{minted}

\subsection*{\code{src/Search/Properties/launchSettings.json}}
\begin{minted}[fontsize=\footnotesize,breakanywhere]{json}
{
  "$schema": "https://json.schemastore.org/launchsettings.json",
  "profiles": {
    "http": {
      "commandName": "Project",
      "dotnetRunMessages": true,
      "launchBrowser": false,
      "applicationUrl": "http://localhost:5004",
      "environmentVariables": {
        "ASPNETCORE_ENVIRONMENT": "Development"
      }
    }
  }
}
\end{minted}

\subsection*{\code{src/Search/Query.cs}}
\begin{minted}{csharp}
using Microsoft.EntityFrameworkCore;

namespace Search;

[QueryType]
public static class Query
{
    [UsePaging(DefaultPageSize = 2, MaxPageSize = 10)]
    public static IQueryable<SessionSearchDocument> SearchSessions(
        SearchContext db,
        double? minimumScore,
        SessionSearchOrder? order)
    {
        var query = db.SessionSearch.AsNoTracking();
        if (minimumScore is not null)
        {
            query = query.Where(
                s => s.AverageScore >= minimumScore.Value);
        }
        return order is SessionSearchOrder.RatingDescending
            ? query.OrderByDescending(s => s.AverageScore)
                .ThenBy(s => s.SessionId)
            : query.OrderBy(s => s.StartsAt)
                .ThenBy(s => s.SessionId);
    }
}
\end{minted}

\subsection*{\code{src/Search/Search.csproj}}
\begin{minted}[fontsize=\footnotesize,breakanywhere]{xml}
<Project Sdk="Microsoft.NET.Sdk.Web">
  <PropertyGroup>
    <TargetFramework>net10.0</TargetFramework>
    <Nullable>enable</Nullable>
    <ImplicitUsings>enable</ImplicitUsings>
    <GenerateDocumentationFile>true</GenerateDocumentationFile>
    <NoWarn>$(NoWarn);CS1591</NoWarn>
  </PropertyGroup>
  <ItemGroup>
    <PackageReference Include="HotChocolate.ApolloFederation"
                      Version="16.6.1" />
    <PackageReference Include="HotChocolate.AspNetCore"
                      Version="16.6.1" />
    <PackageReference Include="HotChocolate.AspNetCore.CommandLine"
                      Version="16.6.1" />
    <PackageReference Include="HotChocolate.Data.EntityFramework"
                      Version="16.6.1" />
    <PackageReference Include="HotChocolate.Types.Analyzers"
                      Version="16.6.1" />
    <PackageReference Include="Microsoft.EntityFrameworkCore.Sqlite"
                      Version="10.0.11" />
  </ItemGroup>
</Project>
\end{minted}

\subsection*{\code{src/Search/SearchContext.cs}}
\begin{minted}{csharp}
using Microsoft.EntityFrameworkCore;

namespace Search;

public sealed class SearchContext(
    DbContextOptions<SearchContext> options) : DbContext(options)
{
    public DbSet<SessionSearchDocument> SessionSearch =>
        Set<SessionSearchDocument>();

    protected override void OnModelCreating(ModelBuilder model)
    {
        model.Entity<SessionSearchDocument>()
            .HasKey(s => s.SessionId);
        model.Entity<SessionSearchDocument>()
            .Property(s => s.StartsAt)
            .HasConversion(
                v => v.UtcDateTime,
                v => new DateTimeOffset(v, TimeSpan.Zero));
        model.Entity<SessionSearchDocument>()
            .HasData(SessionSearchData.Documents);
    }
}
\end{minted}

\subsection*{\code{src/Search/Session.cs}}
\begin{minted}{csharp}
using HotChocolate.Types.Relay;

namespace Search;

public sealed record Session([property: ID] string Id);
\end{minted}

\subsection*{\code{src/Search/SessionSearchData.cs}}
\begin{minted}{csharp}
namespace Search;

public static class SessionSearchData
{
    public static readonly
        IReadOnlyList<SessionSearchDocument> Documents =
    [
        new(1, At(9), 4.0, 3),
        new(2, At(10), 4.5, 2),
        new(3, At(11), 3.0, 1),
        new(4, At(14), null, 0)
    ];

    private static DateTimeOffset At(int hour) =>
        new(2026, 9, 14, hour, 0, 0, TimeSpan.Zero);
}
\end{minted}

\subsection*{\code{src/Search/SessionSearchDocument.cs}}
\begin{minted}{csharp}
namespace Search;

public sealed record SessionSearchDocument(
    [property: GraphQLIgnore] int SessionId,
    DateTimeOffset StartsAt,
    double? AverageScore,
    int RatingCount);
\end{minted}

\subsection*{\code{src/Search/SessionSearchDocumentType.cs}}
\begin{minted}{csharp}
using HotChocolate.Types.Relay;

namespace Search;

[ObjectType<SessionSearchDocument>]
public static partial class SessionSearchDocumentType
{
    // This method cannot return null and the schema says it can, which
    // looks like an error and is the point. What the router puts in
    // this position is not what this method returns: it is whatever
    // the Sessions service answers for the key this method formats. A
    // non-null here is a promise about that service's uptime, made by
    // a service that has no way to keep it, and it costs the whole
    // search page when Sessions is stopped.
    public static Session? GetSession(
        [Parent] SessionSearchDocument document,
        INodeIdSerializer serializer)
        => new(serializer.Format(nameof(Session), document.SessionId));
}
\end{minted}

\subsection*{\code{src/Search/SessionSearchOrder.cs}}
\begin{minted}{csharp}
namespace Search;

public enum SessionSearchOrder
{
    StartsAt,
    RatingDescending
}
\end{minted}

\subsection*{\code{src/Search/SessionType.cs}}
\begin{minted}{csharp}
using HotChocolate.ApolloFederation.Types;

namespace Search;

[ObjectType<Session>]
[Key("id", false)]
public static partial class SessionType;
\end{minted}

\subsection*{\code{src/Search/StatementLog.cs}}
\begin{minted}{csharp}
using System.Data.Common;
using Microsoft.EntityFrameworkCore.Diagnostics;

namespace Search;

public sealed class StatementLog : DbCommandInterceptor
{
    private static int _count;

    public override InterceptionResult<DbDataReader> ReaderExecuting(
        DbCommand command,
        CommandEventData eventData,
        InterceptionResult<DbDataReader> result)
    {
        Console.WriteLine(
            $"-- statement {Interlocked.Increment(ref _count)}");
        Console.WriteLine(command.CommandText);
        return result;
    }

    public override
        ValueTask<InterceptionResult<DbDataReader>> ReaderExecutingAsync(
        DbCommand command,
        CommandEventData eventData,
        InterceptionResult<DbDataReader> result,
        CancellationToken cancellationToken = default)
    {
        Console.WriteLine(
            $"-- statement {Interlocked.Increment(ref _count)}");
        Console.WriteLine(command.CommandText);
        return ValueTask.FromResult(result);
    }
}
\end{minted}

\subsection*{\code{src/Search/appsettings.Development.json}}
\begin{minted}[fontsize=\footnotesize,breakanywhere]{json}
{
  "Logging": {
    "LogLevel": {
      "Default": "Information",
      "Microsoft.AspNetCore": "Warning"
    }
  }
}
\end{minted}

\subsection*{\code{src/Search/appsettings.json}}
\begin{minted}[fontsize=\footnotesize,breakanywhere]{json}
{
  "Logging": {
    "LogLevel": {
      "Default": "Information",
      "Microsoft.AspNetCore": "Warning"
    }
  },
  "AllowedHosts": "*"
}
\end{minted}
\section{Composition and Router Configuration}

\subsection*{\code{graph/graph.yaml}}
\begin{minted}{yaml}
version: 1
subgraphs:
  - name: sessions
    routing_url: http://localhost:5001/graphql
    schema:
      file: ../src/Sessions/schema.graphql
  - name: speakers
    routing_url: http://localhost:5002/graphql
    schema:
      file: ../src/Speakers/schema.graphql
  - name: ratings
    routing_url: http://localhost:5003/graphql
    schema:
      file: ../src/Ratings/schema.graphql
  - name: search
    routing_url: http://localhost:5004/graphql
    schema:
      file: ../src/Search/schema.graphql
\end{minted}

\subsection*{\code{router/config.yaml}}
\begin{minted}{yaml}
version: "1"

# Without this the router asks Cosmo Cloud for its execution config and
# refuses to start, because it has no token to ask with. The path is
# resolved against the directory the router is started in, not against
# this file, so start it from this folder:
#
#   router -config config.yaml
execution_config:
  file:
    path: ../graph/router.json

# A development switch, and what makes the query plan visible at all.
# The plan is an Advanced Request Tracing option, and outside dev mode
# the router wants a Cosmo Cloud token before it hands one out. With
# this off, X-WG-Include-Query-Plan is ignored rather than refused.
dev_mode: true

# Who the router thinks is asking. The alternative entry in this list
# is a jwks url, which is what a deployment with an identity provider
# behind it uses; the shared secret is here because this book stands
# up no such provider. header_key_id is not optional: the router
# matches it against the kid in the token header, so the token this
# book mints carries one.
#
# A request carrying a token this key does not verify is answered 401
# by the router and never reaches a subgraph.
authentication:
  jwt:
    jwks:
      - secret: splitting-the-graph-development-signing-key-32b
        symmetric_algorithm: HS256
        header_key_id: dev

# The router does not forward request headers to subgraphs unless it
# is told to, and a subgraph that guards its own data cannot tell who
# is asking without this block. Leaving it out does not make the graph
# insecure; it makes the guarded field unreachable. Every request that
# names one comes back null with an authorization error underneath,
# for a caller holding a good token exactly as for a caller holding
# none, which is what makes the missing block read like a policy
# rather than like a header. A request that never names a guarded
# field is unaffected, so this survives a smoke test.
headers:
  all:
    request:
      - op: propagate
        named: Authorization
\end{minted}
